\section{Analyse de complexité}
\subsubsection{Notation}
Soit:
\begin{itemize}
    \item $N$: nombre total de termes,
    \item $m$: nombre de termes correspondant au préfixe de la requête,
    \item $k$: nombre de résultats à afficher.
\end{itemize}

\subsection{Approche 1: Recherche linéaire}
Chaque requête parcourt l’ensemble des $N$ termes pour sélectionner ceux correspondant au préfixe,
puis conserve les $k$ meilleurs selon leur poids.

\subsubsection{Complexité en temps}
\textbf{Meilleur cas:}
\[
O(N)
\]
Même dans le meilleur cas, tous les termes doivent être examinés.

\textbf{Pire cas:}
\[
O(N)
\]
Le filtrage est toujours linéaire.

\subsubsection{Complexité totale}
\[
O(N)
\]

\subsubsection{Complexité en espace}
\[
O(1)
\]

\subsection{Approche 2: Tri lexicographique + recherche binaire + tri des résultats}
\begin{itemize}
    \item Les termes sont triés lexicographiquement.
    \item Une recherche binaire permet d'identifier les $m$ termes correspondant au préfixe.
    \item Les $m$ termes sont ensuite triés par poids pour sélectionner les $k$ meilleurs.
\end{itemize}

\subsubsection{Complexité en temps}

\textbf{Tri initial:}
\begin{itemize}
    \item MergeSort: $O(N \log N)$
    \item QuickSort:
    \begin{itemize}
        \item Meilleur cas: $O(N \log N)$
        \item Pire cas: $O(N^2)$
    \end{itemize}
\end{itemize}

\textbf{Recherche binaire:}
\[
O(\log N)
\]
Recherchons 2 éléments extrêmes

\textbf{Extraction des $m$ éléments:}
\[
O(m)
\]

\textbf{Tri des $m$ éléments par poids:}
\begin{itemize}
    \item MergeSort: $O(m \log m)$
    \item QuickSort:
    \begin{itemize}
        \item Meilleur cas: $O(m \log m)$
        \item Pire cas: $O(m^2)$
    \end{itemize}
\end{itemize}

\subsubsection{Complexité totale}

\textbf{Avec MergeSort:}
\[
O(N \log N + m \log m)
\]

\textbf{Avec QuickSort:}
\begin{itemize}
    \item Meilleur cas:
    \[
    O(N \log N + m \log m)
    \]
    \item Pire cas:
    \[
    O(N^2 + m^2)
    \]
\end{itemize}

\subsubsection{Complexité en espace}
\begin{itemize}
    \item MergeSort: $O(N)$
    \item QuickSort:
    \begin{itemize}
        \item Meilleur cas: $O(\log N)$
        \item Pire cas: $O(N)$
    \end{itemize}
\end{itemize}

Stockage supplémentaire pour les résultats:
\[
O(m)
\]

\begin{itemize}
    \item MergeSort: $O(m)$
    \item QuickSort:
    \begin{itemize}
        \item Meilleur cas: $O(\log m)$
        \item Pire cas: $O(m)$
    \end{itemize}
\end{itemize}

\subsubsection{Complexité totale}
\[
O(N)
\]

\subsection{Approche 3: Tri lexicographique + recherche binaire + tas binaire}
\begin{itemize}
    \item Les termes sont triés lexicographiquement.
    \item Une recherche binaire identifie les $m$ termes du préfixe.
    \item Une structure de tas binaire (heap) est utilisée pour sélectionner les $k$ meilleurs selon leur poids.
\end{itemize}

\subsubsection{Complexité en temps}

\textbf{Tri initial:}
\begin{itemize}
    \item MergeSort: $O(N \log N)$
    \item QuickSort:
    \begin{itemize}
        \item Meilleur cas: $O(N \log N)$
        \item Pire cas: $O(N^2)$
    \end{itemize}
\end{itemize}

\textbf{Recherche binaire:}
\[
O(\log N)
\]
Recherchons 2 éléments extrêmes.

\textbf{Parcours des $m$ éléments avec tas de taille $\min(k, m)$:}
\[
O(m \log \min(k, m))
\]

\subsubsection{Complexité totale}

\textbf{Avec MergeSort:}
\[
O(N \log N + m \log \min(k, m))
\]

\textbf{Avec QuickSort:}
\begin{itemize}
    \item Meilleur cas:
    \[
    O(N \log N + m \log \min(k, m))
    \]
    \item Pire cas:
    \[
    O(N^2 +m \log \min(k, m))
    \]
\end{itemize}

\subsubsection{Complexité en espace}
\begin{itemize}
    \item MergeSort: $O(N)$
    \item QuickSort:
    \begin{itemize}
        \item Meilleur cas: $O(\log N)$
        \item Pire cas: $O(N)$
    \end{itemize}
\end{itemize}

Espace supplémentaire pour le tas:
\[
O(\min(k, m))
\]

\subsubsection{Complexité totale}
\[
O(N)
\]
